\subsection{Locally convex Lusin spaces}

PROPOSITION 1. --- Let E be a Hausdorff locally convex space. Suppose that there exists a sequence \((\mathrm{E}_{n})_{n\in\mathbb{N}}\) of Fréchet spaces satisfying the first axiom of countability, and continuous linear mappings \(u_{n}:E_{n}\to E\) such that \(\mathrm{E}=\bigcup_{n\in\mathbb{N}}u_{n}(\mathrm{E}_{n})\) . Then E is a Lusin space.

Let \(P_n\) be the kernel of \(u_n\) ; then \(u_n\) defines a bijective continuous mapping from the quotient space \(E_n / P_n\) onto \(u_n(E_n)\) . Since \(E_n / P_n\) is a Fréchet space satisfying the first axiom of countability (GT, IX, § 3, No.~1), hence a polish space (GT, IX, § 6, No.~1, def. 1), \(u_n(E_n)\) is a Lusin subspace of E (GT, IX, § 6, No.~4, prop. 11). Therefore by GT, IX, § 6, No.~7, cor. of th. 3, the space E, which is regular (GT, III, § 3, No.~1) is a Lusin space.

Example 1. --- Every Fréchet space satisfying the first axiom of countability is a polish space, hence a Lusin space. Consequently, so are the spaces \(\mathcal{C}(\mathrm{X})\) , where \(\mathrm{X}\) is locally compact and has a countable base (the topology of \(\mathcal{C}(\mathrm{X})\) being that of compact convergence, cf.~GT, X, § 3, No.~3, corollary and § 1, No.~6, cor. 3); * the spaces \(\mathcal{C}^{\infty}(\mathrm{U})\) , where U is an open subset of \(\mathbf{R}^{n}\) (III, p.~9) and \(\mathcal{H}(\mathrm{U})\) , where U is an open subset of \(\mathbf{C}^{n}\) (III, p.~10).

Prop. 1 shows that the spaces \(\mathcal{C}_0^\infty (\mathbf{U})\) , where \(\mathbf{U}\) is an open set in \(\mathbf{R}^n\) , \(\mathcal{G}_s(\mathbf{I})\) , where \(\mathbf{I}\) is a compact interval in \(\mathbf{R}\) and \(s\geqslant 1\) , and \(\mathcal{H}(\mathbf{K})\) , where \(\mathbf{K}\) is a compact subset of \(\mathbf{C}^n\) are all Lusin spaces (III, p.~10).

THEOREM 2. --- Let E be a locally convex space, which is the inductive limit of an increasing sequence \((\mathrm{E}_n)_{n\in \mathbb{N}}\) of subspaces of E, endowed with the topologies of Fréchet spaces satisfying the first axiom of countability. Suppose that every compact subset of E is contained in one of the \(E_{n}\) and is compact in this space. Let F be a Fréchet space satisfying the first axiom of countability. Then the space \(\mathcal{L}_{c}(E;F)\) is a Lusin space.

The space E is bornological (III, p.~12), hence the space \(\mathcal{L}_c(\mathrm{E};\mathrm{F})\) is complete (III, p.~23, prop. 12). The linear mapping \(j\colon f\mapsto (f|\mathrm{E}_n)_{n\in \mathbf{N}}\) is an injection from \(\mathcal{L}_c(\mathrm{E};\mathrm{F})\) into the product space \(\prod_{n\in \mathbf{N}}\mathcal{L}_c(\mathrm{E}_n;\mathrm{F})\) ; by virtue of the hypothesis on the compact subsets of E and the definition of the \(\mathfrak{S}\) -topologies, \(j\) is an isomorphism from \(\mathcal{L}_c(\mathrm{E};\mathrm{F})\) onto its image (endowed with the topology induced by the product topology); moreover, since \(\mathcal{L}_c(\mathrm{E};\mathrm{F})\) is complete, this image is a closed subspace of \(\prod_{n\in \mathbf{N}}\mathcal{L}_c(\mathrm{E}_n;\mathrm{F})\)

(GT, II, § 3, No.~4, prop. 8). By GT, IX, § 6, No.~4, it is therefore enough to prove that each of the spaces \(\mathcal{L}_c(\mathrm{E};\mathrm{F})\) is a Lusin space. For the rest of the proof, we shall assume that E is a Fréchet space satisfying the first axiom of countability.

Since F is a Fréchet space satisfying the first axiom of countability, it is isomorphic to a closed subspace of a countable product of Banach spaces \(F_{n}\) , each of which is a quotient of F (II, p.~5), hence satisfies the first axiom of countability. The linear mapping \(j': f \mapsto (\mathrm{pr}_{n} \circ f)_{n \in \mathbb{N}}\) is an injection from \(\mathcal{L}_{c}(\mathrm{E}; \mathrm{F})\) into the product space \(\prod_{n \in \mathbb{N}} \mathcal{L}_{c}(\mathrm{E}; \mathrm{F}_{n})\) , and by using the definition of the S-topologies and of the open sets in a product, \(j'\) is an isomorphism from \(\mathcal{L}_{c}(\mathrm{E}; \mathrm{F})\) onto its image; moreover, since \(\mathcal{L}_{c}(\mathrm{E}; \mathrm{F})\) is complete, this image is a closed subspace of \(\prod_{n \in \mathbb{N}} \mathcal{L}_{c}(\mathrm{E}; \mathrm{F})\) . Therefore it is enough to prove that each of the spaces \(\mathcal{L}_{c}(\mathrm{E}; \mathrm{F}_{n})\) is a Lusin space (GT, IX, § 6, No.~4), and consequently, we can assume that F is a Banach space satisfying the first axiom of countability.

The space \(\mathcal{L}_{c}(\mathrm{E};\mathrm{F})\) is the union of a countable family of equicontinuous and closed subsets (III, p.~19, cor. 1 and GT, X, § 2, No.~3, prop. 6). But every equicontinuous subset H of \(\mathcal{L}_{c}(\mathrm{E};\mathrm{F})\) is metrizable and satisfies the first axiom of countability (III, p.~18, prop. 6 and GT, X, § 2, No.~4, th. 1); if H is closed, then it is a complete space for the uniform structure induced by that of \(\mathcal{L}_{c}(\mathrm{E};\mathrm{F})\) , since the latter is complete In other words, H is a polish space, and a fortiori a Lusin space; consequently the regular space \(\mathcal{L}_{c}(\mathrm{E};\mathrm{F})\) is a Lusin space (GT, IX, § 6, No.~7, cor. of th. 3).

COROLLARY. --- The hypotheses on E being as in th. 2, assume, in addition that every bounded subset of E is relatively compact. Then the strong dual of E is a Lusin space. * In particular, the strong dual of a Fréchet space satisfying the first axiom of countability, which is also a Montel space, is a Lusin space. *

* Example 2. --- Let U be an open subset of \(\mathbf{R}^{n}\) . The corollary applies in particular to the Fréchet space \(E = \mathcal{C}^{\infty}(\mathrm{U})\) ; its dual \(\mathcal{C}_{0}^{-\infty}(\mathrm{U})\) (the space of distributions with compact support on U) is then a Lusin space. The space \(\mathcal{C}_{0}^{\infty}(\mathrm{U})\) is a strict inductive limit of a sequence of Fréchet spaces \(\mathcal{C}_{K_n}^{\infty}(\mathrm{U})\) satisfying the first axiom of countability (III, p.~9). We can show that each of the spaces \(\mathcal{C}_{K_n}^{\infty}(\mathrm{U})\) is a Montel space; in addition, every bounded subset of \(\mathcal{C}_{0}^{\infty}(\mathrm{U})\) is contained in one of the spaces \(\mathcal{C}_{K_n}^{\infty}(\mathrm{U})\) (III, p.~5, prop. 6). We can then apply the corollary of th. 2. Then the dual \(\mathcal{C}_{0}^{-\infty}(\mathrm{U})\) of \(\mathcal{C}_{0}^{\infty}(\mathrm{U})\) (the space of distributions on U) is a Lusin space for the strong topology. *

Similarly we prove that for every open subset U of \(C^{n}\) , and for every compact subset K of \(C^{n}\) , the strong dual of \(\mathcal{H}(U)\) and the strong dual of \(\mathcal{H}(K)\) are Lusin spaces. *

Remark. --- Let E be as in th. 2; let F be a Hausdorff locally convex space which is the union of the images of a sequence of continuous linear mappings \(u_n: \mathrm{F}_n \to \mathrm{F}\) , where each \(\mathrm{F}_n\) is a Fréchet space satisfying the first axiom of countability; then \(\mathcal{L}_c(\mathrm{E};\mathrm{F})\) is a Lusin space. As in prop. 1, we first reduce to the case where each \(u_n\) is injective; then, as in the proof of th. 2, we can assume that E is a Fréchet space satisfying the first axiom of countability. Then, by I, p.~20, prop. 1, \(\mathcal{L}(\mathrm{E};\mathrm{F})\) is the union of the \(\mathcal{L}(\mathrm{E};\mathrm{F}_n)\) ; moreover, the canonical injection \(\mathcal{L}_c(\mathrm{E};\mathrm{F}_n) \to \mathcal{L}_c(\mathrm{E};\mathrm{F})\) is continuous (GT, X, § 1, No.~4, prop. 3). Since each of the spaces \(\mathcal{L}_c(\mathrm{E};\mathrm{F}_n)\) is a Lusin space by th. 2, \(\mathcal{L}(\mathrm{E};\mathrm{F}_n)\) is also a Lusin space for the topology induced by that of \(\mathcal{L}_c(\mathrm{E};\mathrm{F})\) (GT, IX, § 6, No.~4, prop. 11); consequently \(\mathcal{L}_c(\mathrm{E};\mathrm{F})\) is a Lusin space by virtue of GT, IX, § 6, No.~7, corollary of th. 3.

